\cleardoublepage
\chapter{RINGKASAN BUKTI KEGIATAN YANG TELAH DISANITASI}
Bukti kegiatan ditempatkan pada lampiran agar isi utama laporan tetap fokus pada proses dan hasil pekerjaan. Dokumen internal yang mengandung data perusahaan tidak direproduksi dalam laporan ini.

\begin{table}[htbp]
\caption{Ringkasan bukti kegiatan Kerja Praktik}
\label{tab:bukti-kegiatan}
\centering
\begin{tabularx}{\textwidth}{|p{0.08\textwidth}|p{0.38\textwidth}|X|}
\hline
\textbf{No.} & \textbf{Kegiatan} & \textbf{Bukti yang tersedia} \\
\hline
1 & Pemahaman alur data, campaign, dan leads & Catatan kerja dan diskusi internal \\
\hline
2 & Pengembangan query dan validasi data & Query kerja, hasil pengujian, dan catatan waterfall pada lingkungan internal \\
\hline
3 & Integrasi dan pengujian ETL & Job Talend dan hasil run pada lingkungan internal \\
\hline
4 & Analisis campaign financing & Workbook analisis dan deck presentasi internal \\
\hline
5 & Spatial clustering data log aplikasi & Kode analisis, workbook hasil, dan deck presentasi internal \\
\hline
\end{tabularx}
\end{table}

Tabel \ref{tab:bukti-kegiatan} menunjukkan jenis bukti yang tersedia. File asli tidak dimasukkan ke repository laporan karena sebagian memuat data, struktur, atau output internal. Jika perusahaan mengizinkan bukti visual untuk versi final, screenshot yang telah disanitasi dapat ditambahkan pada lampiran ini tanpa memindahkan data mentah.
